%%%%%%%%%%%%%%%%%%%%%%%%%%%%%%%%%%%%%%%%%%%%%%%%%%%%%%%%%%%%%%%%%%%%%%%%%%%%%%%
%% Resumo --- ABNT NBR 6028:2021
%%   parágrafo único; voz ativa; sem citações;
%%   palavras-chave separadas por ponto e vírgula, iniciais minúsculas,
%%   finalizadas por ponto.
%%
%% Enuncia a MESMA afirmação central do Capítulo 1 e do Capítulo 5, com a mesma
%% força. Ao alterar a tese, alterar os três.
%%%%%%%%%%%%%%%%%%%%%%%%%%%%%%%%%%%%%%%%%%%%%%%%%%%%%%%%%%%%%%%%%%%%%%%%%%%%%%%

\begin{resumo}
A manutenção preditiva de máquinas rotativas por análise de vibração é normalizada, e o
obstáculo à sua adoção é o custo de instrumentação, que restringe a cobertura aos ativos
críticos. Nós sensores de baixo custo respondem a esse obstáculo, mas a literatura que os
propõe raramente declara a banda e o piso de ruído que sustentam o desempenho reportado.
Este trabalho determina o que a cadeia de medição especificada para um dispositivo dessa
classe --- acelerômetro de consumo com taxa de amostragem de 1~kHz, microcontrolador de
baixo consumo, enlace de longo alcance e invólucro impresso em polímero --- permite medir,
e calibra o estágio de decisão de modo que o desempenho seja verificável no ponto de
operação pretendido. O método constrói condições de referência cujo resultado correto
decorre de solução analítica ou de definição normativa, submete-as à cadeia proposta e mede
o desvio. Os resultados têm três procedências, declaradas individualmente: sinal
sintetizado, dado real de bancada decimado em \emph{software} e modelagem sobre as
especificações de projeto. Em dado real, a 2\% de falso alarme, a detecção é de
75,0\% $\pm$ 3,5\% em desbalanceamento, entre 14,8\% e 20,7\% em defeito de rolamento puro
e entre 5,0\% e 7,2\% em desalinhamento. Em sinal sintetizado, três decisões de estimação
alteram os resultados por margens maiores que o efeito físico investigado: a integração no
domínio do tempo erra o indicador normativo em 73\% sem viés e em 716\% sob viés contido na
especificação do sensor, enquanto a integração no domínio da frequência o recupera
exatamente; a política de limiar padrão do detector marca 42,1\% da operação normal como
anomalia, contra 1,4\% sob limiar calibrado; e a partição aleatória de janelas infla a
acurácia em cerca de 36 pontos percentuais frente à partição por ensaio. Por modelagem, a
frequência de montagem do conjunto situa-se em 557~Hz, dentro da banda de medição, e o
limite térmico de operação é imposto pela bateria, e não pelo invólucro. Conclui-se que a
cadeia sustenta a detecção de desbalanceamento sob falso alarme fixado e não sustenta
diagnóstico de defeito de rolamento, e que a contribuição verificável de um dispositivo
dessa classe está no estágio de decisão, desde que o estágio de aquisição esteja declarado.
As conclusões valem para a cadeia especificada, e não para o desempenho do dispositivo
montado em operação industrial: a construção do protótipo e seu ensaio em ambiente real
constituem a etapa seguinte.

\vspace{\onelineskip}
\noindent
\textbf{Palavras-chave}: manutenção preditiva; análise de vibração; detecção de anomalia;
sistemas embarcados; instrumentação de baixo custo.
\end{resumo}
